% Source: https://artofproblemsolving.com/community/c3h1915246_meow_nya_logic_puzzle_2
% Stable identity: aops-topic-1915246-q01
\textbf{A key lemma.} For a yes-or-no question $Q$, let $E(Q)$ be the question ``If I asked you $Q$, would you answer Meow?'' If Truecat or Liecat is asked $E(Q)$, the answer is Meow exactly when $Q$ is true.

\emph{Proof.} Check the four cases.
\begin{itemize}
\item Truecat, Meow $=$ yes: Truecat would answer $Q$ with Meow exactly when $Q$ is true, so $E(Q)$ is true exactly when $Q$ is true, and the truthful answer is Meow ($=$ yes) exactly then.
\item Truecat, Meow $=$ no: Truecat would answer $Q$ with Meow exactly when $Q$ is false; so $E(Q)$ is false exactly when $Q$ is true, and the truthful answer ``no'' is Meow exactly then.
\item Liecat, Meow $=$ yes: Liecat would answer $Q$ with Meow ($=$ yes) exactly when $Q$ is false; so $E(Q)$ is false exactly when $Q$ is true, and the lying answer ``yes'' $=$ Meow comes exactly then.
\item Liecat, Meow $=$ no: Liecat would answer $Q$ with ``no'' $=$ Meow exactly when $Q$ is true; so $E(Q)$ is true exactly when $Q$ is true, and the lying answer ``no'' $=$ Meow comes exactly then.
\end{itemize}
In every case the answer is Meow exactly when $Q$ is true. $\blacksquare$ (The question $E(Q)$ is about facts and the cat's own behaviour, so it causes no paradox.)

\textbf{(a) Three questions.} Call the cats $A$, $B$, $C$.
\begin{enumerate}
\item Ask $A$: $E(\text{``Is }B\text{ Randomcat?''})$. If the answer is Meow, then either $A$ is Randomcat, or $A$ is not and (by the lemma) $B$ is Randomcat; in both cases $C$ is \emph{not} Randomcat. If the answer is Nya, then either $A$ is Randomcat, or $A$ is not and $B$ is not Randomcat; in both cases $B$ is not Randomcat. Either way we know a cat $X\ne A$ that is not Randomcat.
\item Ask $X$: $E(\text{``Are you Truecat?''})$. By the lemma, Meow means $X$ is Truecat and Nya means $X$ is Liecat.
\item Ask $X$: $E(\text{``Is }A\text{ Randomcat?''})$. By the lemma this tells us whether $A$ is Randomcat. The other cat besides $X$ and $A$ then gets the remaining name.
\end{enumerate}
After three questions every cat is identified.

\textbf{(b) Two questions are not enough.} There are $3!=6$ possible ways to assign the names. With two questions there are at most $4$ possible answer sequences (Meow-Meow, Meow-Nya, Nya-Meow, Nya-Nya), even if the second question depends on the first answer. The final conclusion depends only on the answers, so a two-question strategy can only ever announce at most $4$ different assignments. Some assignment is then never announced, and if it is the true one, the strategy fails, whatever the cats answer.

So the least number of questions that always suffices is $\boxed{3}$.
